\section{Method}\label{sec:method}

\scl{} is a small step that runs once per frame between the detector and
the tracker (Fig.~\ref{fig:arch}). In plain terms, it does four things.
First, it looks at the detector scores of the last few frames and sorts the
candidates into three groups: background, ambiguous, and confident
(Sect.~\ref{sec:calib}). Second, it uses these groups to decide whether the
detector's output is clean or noisy (Sect.~\ref{sec:regimes}). Third, it
acts on that decision: on a clean stream it leaves the tracker's own settings
alone, and on a noisy stream it chooses which candidates reach the tracker,
gives them new scores, and sets the tracker's thresholds. Fourth, in both
cases it can pass a weak candidate that clearly continues an existing track
(Sect.~\ref{sec:rescue}). The tracker then runs as usual on what it receives.
The detector and the tracker are never modified, and
Table~\ref{tab:constants} lists every constant of the layer.

\begin{figure*}[t]
\centering
\begin{tikzpicture}[
  box/.style={draw, rounded corners=2pt, align=center, font=\footnotesize, minimum height=9mm, minimum width=22mm, inner sep=3pt},
  arr/.style={-{Stealth[length=2mm]}, thick}, lab/.style={font=\scriptsize}]
\node[box] (det) at (0,0) {Frozen\\detector};
\node[box] (pol) at (4.3,0) {Candidate policy\\(clean or noisy)};
\node[box] (res) at (8.5,0) {Continuation\\rescue};
\node[box] (trk) at (12.7,0) {Frozen\\tracker};
\node[box] (stat) at (0,-2.1) {Stream statistics:\\nested Otsu on logits\\of frames $t-W..t-1$};
\node[box] (reg) at (4.3,-2.1) {Regime estimate\\clean if $\bar\rho_t \ge \tfrac12$};
\node[box] (hc) at (10.6,-2.1) {Host contract\\$(a, b, \ell, m_0)$};
\draw[arr] (det) -- node[lab, above]{candidates} (pol);
\draw[arr] (pol) -- (res);
\draw[arr] (res) -- node[lab, above]{$K_t$, scores} (trk);
\draw[arr] (det) -- (stat);
\draw[arr] (stat) -- node[lab, above]{$t_1, t_2, \rho$} (reg);
\draw[arr] (reg) -- (pol);
\draw[arr] (hc.north) ++(-0.6,0) |- ++(0,0.45) -| (res.south);
\draw[arr] (hc.north) ++(0.6,0) |- ++(0,0.45) -| node[lab, pos=0.75, right]{$a_t, b_t, m_t$} (trk.south);
\draw[arr] (trk.north) -- ++(0,0.45) -| node[lab, pos=0.25, above]{tracks of frame $t-1$} (res.north);
\end{tikzpicture}
\caption{Self-calibrating control layer (\scl{}) between a frozen detector and a
frozen tracker. The statistics use earlier frames only. The host contract is the
tracker's published operating point. The tracker receives the selected
candidates, their scores, and its thresholds for the frame.}
\label{fig:arch}
\end{figure*}

\subsection{Problem setting}\label{sec:problem}
In each frame $t$, the detector gives a list of candidates
$D_t=\{(\mathbf{x}_i,s_i,c_i)\}$: a box $\mathbf{x}_i$, a score
$s_i\in(0,1)$, and a class label $c_i$ for each one, down to a lowest score
called the emission floor. The layer does not use the score directly but its
logit (log-odds), which spreads out scores close to 0 and 1:
\begin{equation}
z_i=\log\frac{s_i}{1-s_i}.
\label{eq:logit}
\end{equation}
The tracker, called the host, tells the layer four numbers about itself,
the host contract $h=(a,b,\ell,m_0)$. These are simply the settings that its
authors published: $a$ is the score a detection needs to be matched to a
track in the first step, $b$ is the score it needs to start a new track,
$\ell$ is the lowest score the tracker uses at all, and $m_0$ is its matching
tolerance (one minus the smallest overlap that still counts as a match). A
single-stage tracker uses one threshold, so $\ell=\min(a,b)$. A two-stage
tracker also uses weaker detections in a second step, so
$\ell<\min(a,b)$.

In every frame the layer returns five things: the candidates $K_t$ that
the tracker will see, their scores $\tilde{s}$, and the thresholds
$(a_t,b_t)$ and tolerance $m_t$ that the tracker will use in this frame:
\begin{equation}
\left(K_t,\tilde{s},a_t,b_t,m_t\right)=
\mathcal{L}\left(Z_{<t},D_t,\mathcal{T}_{t-1},g_t;\,h\right).
\label{eq:layer}
\end{equation}
Here $Z_{<t}$ are the logits of the earlier frames, $\mathcal{T}_{t-1}$
are the tracks of the previous frame, and $g_t$ is an optional measure of
camera motion in frame $t$. The decision for frame $t$ uses only past
frames; frame $t$ is added to the statistics after the decision, and no
future frame is ever used.

Four requirements guided the design. The layer needs no training and no
labels. One frozen configuration must serve every tracker, detector, and
dataset. When the tracker already fits the detector, its input must stay
unchanged. And the layer must cost little time next to the detector.


\subsection{Self-calibration from the score stream}\label{sec:calib}
A detector that suits the scene gives many low scores to background and a
separate group of high scores to objects. The layer checks how clearly this
happens, without labels. It collects the logits of all candidates of the
last $W=10$ frames and splits them twice with Otsu's method \cite{otsu1979},
which finds the threshold that best separates a histogram into two groups.
The first split, $t_1$, separates background from foreground. The second
split, $t_2$, divides the foreground into ambiguous and confident
candidates. The share of confident candidates in the foreground,
\begin{equation}
\rho_t=\frac{\left|\{i: z_i\ge t_2\}\right|}{\left|\{i: z_i\ge t_1\}\right|},
\label{eq:rho}
\end{equation}
tells how cleanly the detector separates objects from clutter. A high
$\rho_t$ means that most foreground candidates are confident; a low
$\rho_t$ means that many of them are ambiguous. Because the splits are
made on logits, they move together with the scores when the whole logit
scale is stretched or shifted, for example by a change of temperature, and
each group keeps the same candidates.

These groups only make sense if there is a real background group. Some
detectors give very few low scores; then the first split falls inside the
objects. The layer detects this case, when fewer candidates lie below $t_1$
than above it, and counts the frame as clean, $\rho_t=1$. This background
check was added after a development case: at the start of some sequences,
with about ten candidates per frame, the splits fell at $t_1\approx0.6$ and
$t_2\approx0.93$ in probability terms and gave a wrong noisy decision.

\subsection{Intervention regimes}\label{sec:regimes}
The regime decides how much the layer is allowed to change. A clean stream
means that the detector already separates objects from clutter, so the
tracker's own settings are kept. A noisy stream means that many candidates
are ambiguous, so the layer takes over the choice of candidates and
thresholds.

\emph{Regime estimate.} The stream is clean at frame $t$ if the median
$\bar\rho_t$ of all $\rho$ values so far (ignoring cold frames) is at
least one half, and noisy otherwise. Before the first valid split, for
example in the first frame, the regime is called cold. Because the median
covers the whole stream so far, the regime describes the detector and the
scene, and a short drop in confidence, for example during fast camera
motion, does not change it.

\emph{Cold and clean regimes.} The tracker keeps its own thresholds. In
the clean regime there is one exception: if a threshold is so high that it
would reject the confident group, it is lowered to $t_2$:
\begin{equation}
a_t=\min\!\left(a,\sigma(t_2)\right),\qquad b_t=\min\!\left(b,\sigma(t_2)\right),
\label{eq:clean}
\end{equation}
where $\sigma$ is the logistic function, which turns a logit back into a
score. Scores are passed unchanged, and only cross-class duplicates (one box
proposed with two class labels) are removed. When the confident group
already lies above the tracker's thresholds, the tracker therefore receives
exactly its normal input.

\emph{Example.} When the detector scores are halved, no score exceeds 0.5, so
a tracker with $a=0.6$ and $b=0.7$, such as ByteTrack in its official setting
(Table~\ref{tab:contracts}), receives no detection that can start a track,
so no track is ever created. The Otsu splits are computed from the halved scores themselves, so
they move with them. In the clean regime, Eq.~\eqref{eq:clean} then lowers
the tracker's thresholds to $\sigma(t_2)$, and the confident candidates reach
the tracker again. In the noisy regime, the remap below places them at or
above the thresholds of 0.5. Section~\ref{sec:shift} shows the measured effect.

\emph{Noisy regime.} Here the layer decides which candidates the tracker
may use, and for what. Confident candidates ($z\ge t_2$) may start and
continue tracks. Ambiguous candidates ($t_1\le z<t_2$) may only continue
existing tracks. Background candidates ($z<t_1$) are dropped. Each score is
replaced by a new one based on its rank $u_i$ among the recent scores of the
stream (from 0 for the lowest to 1 for the highest),
\begin{equation}
\tilde{s}_i=
\begin{cases}
0.5+0.5\,u_i, & z_i\ge t_2,\\
0.1+0.4\,u_i, & t_1\le z_i<t_2,
\end{cases}
\label{eq:remap}
\end{equation}
and the tracker's matching and birth thresholds are both set to 0.5,
$a_t=b_t=0.5$. Confident candidates thus get scores of 0.5 and above, and
ambiguous ones get scores below 0.5, so they can only reach a tracker's
low-score step. Duplicates are removed with the help of the tracks: a weaker
candidate that overlaps a stronger one with an intersection over union (IoU)
above one half is removed, unless a different track of the previous frame
claims it. This keeps people who stand close together in a crowd.

\subsection{Continuation rescue and motion rule}\label{sec:rescue}
A tracker ignores every candidate below its lowest threshold, so a track
can be lost when its object scores low for a few frames. This matters most
for single-stage trackers, which have no low-score step. The rescue gives
such a candidate a second chance, but only when it clearly continues an
existing track. Three conditions must hold: the candidate is in the
foreground ($z\ge t_1$, or any candidate for a single-stage tracker when
there is no background group); it overlaps a track of the previous frame
with an IoU of at least one half; and no other usable candidate covers that
track. In addition, the tracker must be unable to see the candidate with the
thresholds of this frame. The candidate is then passed with its score raised just above
the tracker's lowest threshold. For a two-stage tracker, this only concerns
candidates below its low threshold $\ell$. For a single-stage tracker, it
concerns candidates below the lowest threshold passed for this frame, so the rescue gives it the
continuation step it lacks.

When the camera moves fast, a box moves more between two frames, and a
fixed overlap requirement can break a correct match. In noisy frames, and
only when a motion measure is available, the matching tolerance is
therefore widened as the current camera motion grows; $r_t$ is the current
motion divided by its median over past frames. The motion is the shift between two
consecutive grayscale frames at quarter resolution, found by phase
correlation \cite{reddy1996phasecorr} and divided by the image diagonal.
The tolerance is
\begin{equation}
m_t=\min\!\left(0.95,\;1-\frac{1-m_0}{\max(1,r_t)}\right).
\label{eq:motion}
\end{equation}

\subsection{Host contract and adapters}\label{sec:contract}
Table~\ref{tab:contracts} lists the contracts used in this paper. Each
value is taken from the tracker's published configuration; none is fitted.
A small adapter for each tracker family passes the layer's decision to the
tracker: the selected candidates and their scores replace the detector
output of the frame, and the new thresholds and tolerance replace the
tracker's own for that frame. Any post-processing that the authors apply
after tracking, such as interpolation or offline linking, is kept. The
policy and its constants are the same for every tracker; a tracker enters
only through its four contract values, and the adapter only translates the
decision, it never makes one.

\begin{table*}[t]
\caption{Host contracts: the operating point that each published tracker
declares, as used by the layer. $\tau$ is the tracker's own published
threshold, which differs between trackers.}
\label{tab:contracts}
\centering\footnotesize
\setlength{\tabcolsep}{4pt}
\begin{tabular}{@{}>{\raggedright\arraybackslash}p{30mm}lcccc>{\raggedright\arraybackslash}p{32mm}@{}}
\toprule
Tracker (setting) & Type & assoc $a$ & birth $b$ & low $\ell$ & match $m_0$ & Decision mapped to\\
\midrule
ByteTrack (official MOT17) & two-stage & 0.6 & 0.7 & 0.1 & 0.8 & thresholds, candidates, scores\\
ByteTrack (library default) & two-stage & 0.25 & 0.25 & 0.1 & 0.8 & same\\
SparseTrack & two-stage & $\tau$ & $\tau+0.1$ & 0.1 & published & same\\
BoostTrack & two-stage & $\tau$ & $\tau$ & 0.1 & $1-$IoU gate & same\\
Hybrid-SORT & two-stage & $\tau$ & $\tau$ & 0.1 & $1-$IoU gate & same\\
OC-SORT & single-stage & 0.6 & 0.6 & 0.6 & 0.7 & same\\
PD-SORT & single-stage & $\tau$ & $\tau$ & $\tau$ & $1-$IoU gate & same\\
\bottomrule
\end{tabular}
\end{table*}


\subsection{Frozen constants}\label{sec:constants}
Table~\ref{tab:constants} lists every constant of the layer; the same
values are used for every tracker, detector, and dataset. The thresholds
$t_1$ and $t_2$ are not constants: they are recomputed in every frame. Each
constant has a simple origin. The overlap of one half for duplicates and
rescue follows MOT evaluation, where two boxes with IoU above 0.5 cannot both
match different objects. The score ranges of Eq.~\eqref{eq:remap} put the
two groups on the common thresholds of 0.5 (matching and birth) and 0.1 (low
score). The cap of 0.95 keeps a minimum overlap of 0.05 for every match. The
window, the motion history, and the remap memory are default values that
were kept fixed in all development stages.

\begin{table}[t]
\caption{Frozen constants of \scl{}.}
\label{tab:constants}
\centering\footnotesize
\begin{tabular}{@{}>{\raggedright\arraybackslash}p{39mm}>{\raggedright\arraybackslash}p{31mm}@{}}
\toprule
Constant & Value\\
\midrule
Statistics window $W$ & 10 frames\\
Regime rule & clean if median $\rho\ge0.5$ over non-cold frames so far\\
Background check & $\rho=1$ if fewer candidates below $t_1$ than above\\
Noisy score ranges & primary $[0.5,1)$, extension $[0.1,0.5)$\\
Noisy association and birth thresholds & 0.5\\
Remap memory & 20 frames, one every 10 frames\\
Duplicate and rescue IoU & 0.5\\
Rescued score & lowest used threshold $+\,0.001$\\
Motion cue & phase correlation, quarter resolution\\
Motion history, warm-up & 100, 5 frames\\
Tolerance cap & 0.95\\
\bottomrule
\end{tabular}
\end{table}

\subsection{Per-frame procedure and cost}\label{sec:cost}
Algorithm~\ref{alg:layer} gives the steps for one frame. The work per
frame is small: the layer sorts the logits of at most $W$ frames for the two
Otsu splits, computes overlaps between candidates and with the tracks of the
previous frame, and, when motion is used, computes one phase correlation on
a quarter-resolution image. With $n$ candidates per frame and $k$ tracks,
the cost is $O(Wn\log(Wn)+n^2+nk)$ plus the fixed image cost. The layer
stores only the logits of the last $W$ frames, the per-frame $\rho$ values,
the remap memory, and the last 100 motion values.

\begin{algorithm}[t]
\caption{\scl{} for frame $t$}\label{alg:layer}
\begin{algorithmic}[1]
\Require candidates $D_t$, tracks $\mathcal{T}_{t-1}$, cue $g_t$, contract $(a,b,\ell,m_0)$
\State $z_i\gets$ logit of $s_i$ for every candidate
\If{no valid split exists}
  \State regime $\gets$ cold; $(a_t,b_t)\gets(a,b)$
\Else
  \State $(t_1,t_2)\gets$ nested Otsu on the window
  \State $\rho_t\gets$ Eq.~\eqref{eq:rho} with the background check
  \State regime $\gets$ clean if $\mathrm{median}(\rho)\ge0.5$, else noisy
\EndIf
\If{regime is cold or clean}
  \State remove cross-class duplicates; keep raw scores
  \If{clean} set $(a_t,b_t)$ by Eq.~\eqref{eq:clean}
  \EndIf
\Else
  \State remove duplicates with track context
  \State drop $z<t_1$; remap scores by Eq.~\eqref{eq:remap}
  \State $a_t\gets0.5$; $b_t\gets0.5$
\EndIf
\State rescue track continuations (not in cold frames)
\State $m_t\gets$ Eq.~\eqref{eq:motion} if noisy, else $m_0$
\State pass $K_t$, $\tilde{s}$, $a_t$, $b_t$, $m_t$ to the tracker
\State add the logits of frame $t$ to the window; add $g_t$ to the history
\State store the tracks of frame $t$ after the tracker update
\end{algorithmic}
\end{algorithm}
